\documentclass[11pt]{article}
\usepackage[T1]{fontenc}
\usepackage{lmodern,amsmath,amssymb,booktabs}
\usepackage[margin=1in]{geometry}
\newcommand{\bin}[2]{\binom{#1}{#2}}
\begin{document}
\begin{center}
{\Large\bfseries How many copies of each color?}\\[4pt]
{\large Rainbow perfect matchings in a randomly colored even cycle}\\[5pt]
Discrete Math Lab \quad$\cdot$\quad Working note
\end{center}

Fix $q$ colors and choose their positive multiplicities before shuffling the
color tokens uniformly over the edges of $C_{2n}$, $n\ge2$. Which
multiplicities maximize the probability of a rainbow perfect matching?
Already on a square, two red and two blue tokens succeed with probability
$2/3$, while three red and one blue always succeed.

\paragraph{Classification.}
If $q<n$, success is impossible. For $n\le q\le 2n$ put
$d=2n-q$ and $\ell=q-n$. The balanced inventory $B$ consists of
$d$ doubles and $2\ell$ singles. When $d\ge2$, let $T$ be obtained from
$B$ by replacing two doubles with a triple and a single. The table lists
all maximizing inventories up to color names.
\begin{center}
\begin{tabular}{lll}
\toprule
Doubles in $B$ & Range & Maximizer\\
\midrule
$0,1$ & all & $B$\\
$2,3$ & all & $T$\\
$4$ & $\ell\le6$ / $\ell\ge7$ & $T$ / $B$\\
$5$ & $\ell\le1$ / $\ell\ge2$ & $T$ / $B$\\
$6$ & $\ell=0$ / $\ell=1$ / $\ell\ge2$ & $T$ / both / $B$\\
$\ge7$ & $\ell=0$ / $\ell\ge1$ & $T$ / $B$\\
\bottomrule
\end{tabular}
\end{center}

\paragraph{Counting successful assignments.}
Distinguish equally colored tokens. One alternating matching receives a
uniform $n$-subset of the $2n$ tokens, and the other receives its complement.
For a multiplicity vector $c$, the elementary symmetric sum $e_n(c)$ counts
subsets with distinct colors. By inclusion--exclusion the success probability is
\[
\frac{2e_n(c)-J(c)}{\bin{2n}{n}},
\qquad
J(c)=
\begin{cases}
0 & \max c_i\ge3,\\
2^d\bin{2\ell}{\ell} & c=B .
\end{cases}
\]
Indeed, for both halves to be rainbow each double must split, while half
the singles go to each side. A triple cannot split between two rainbow halves.

\paragraph{Reduction to two candidates.}
Among inventories with a triple or larger entry, maximize $e_n(c)$.
Moving one unit from $a$ to $b$ with $a\ge b+2$ gains
$(a-b-1)e_{n-2}$ of the remaining entries, strictly positive because
$q\ge n$. Reduce a largest entry to three by transfers to singles, which
exist whenever the mean is at most two and a larger entry occurs.
Keep that triple fixed and balance the rest. Their mean is below two, so
the result is $T$. Every changed inventory strictly improves along this path.

\newpage
\paragraph{Comparing the two candidates.}
Assume $d\ge2$. Moving from $T$ to $B$ increases $e_n$ by $E$, the
elementary symmetric sum of degree $n-2$ on $d-2$ doubles and $2\ell$
singles. It also incurs the intersection penalty $J(B)$. Define
\[
R=\frac{E}{2^{d-2}\bin{2\ell}{\ell}}
=\sum_{i=0}^{\min(d-2,\ell)}
\bin{d-2}{i}2^{-i}
\frac{\bin{2\ell}{\ell+i}}{\bin{2\ell}{\ell}} .
\]
Here $i$ counts omitted doubles and $\ell+i$ singles must be selected.
Consequently $2E-J(B)$ has the sign of $R-2$. A positive sign favors $B$.

For $\ell=0$, $R=1$. For $\ell\ge1$, $R$ is nondecreasing in $d$ and $\ell$:
in the binomial ratio, each factor $(\ell-j)/(\ell+j+1)$ for
$0\le j<i$ increases with $\ell$, with a zero term until $\ell\ge i$.
For $d=2,3$ the sum stays below two. For $d=4$ direct simplification gives
\[
R-2=\frac{\ell(\ell-1)-4(\ell+2)}{4(\ell+1)(\ell+2)},
\]
which changes sign between six and seven. For $d=5$ one has
$R(5,1)=7/4$ and $R(5,2)=17/8$. Finally
$R(d,1)=1+(d-2)/4$, and the term for $i=1$ increases strictly with
$\ell$ for $d\ge3$. These comparisons give the table, including its
single tie at $(n,q)=(7,8)$.

\paragraph{Size of the effect.}
When $q=n$, both candidates admit closed forms. The balanced inventory
has success probability $2^n/\bin{2n}{n}$, while the optimal inventory
has probability $6\cdot2^{n-2}/\bin{2n}{n}$. For every $n$, introducing
one triple therefore increases success by a factor of $3/2$.

\paragraph{Verification.}
The accompanying program \texttt{check\_cycle\_inventory.py} enumerates
individual token subsets for every multiplicity partition through twelve
edges. It also compares the table with exact counts for every partition
with $n\le q\le2n$ through thirty edges. These finite checks supplement
the proof. Priority has not been established.
\end{document}
